% TOP_PROBLEM: {"problemId": "problem.log-rank-conjecture-communication-complexity", "canonicalTitle": "Log-Rank Conjecture in Communication Complexity", "releaseRank": 144, "primaryDomain": "cryptography_coding_information_optimization", "primaryDomainLabel": "Cryptography, coding, information & optimization", "displayStatus": "open", "releaseStatus": "open"}
% !TeX program = lualatex
% Definition and sources only; this file is not an LLM solution attempt.
\documentclass[11pt]{article}
\usepackage[margin=25mm]{geometry}
\usepackage{fontspec}
\setmainfont{Latin Modern Roman}
\usepackage{amsmath,amssymb,mathtools}
\usepackage{unicode-math}
\setmathfont{Latin Modern Math}
\usepackage{xurl,hyperref}
\hypersetup{colorlinks=true,urlcolor=blue}
\setcounter{secnumdepth}{0}
\setlength{\parindent}{0pt}
\setlength{\parskip}{5pt}
\setlength{\emergencystretch}{3em}
\begin{document}
\section*{\texorpdfstring{Log-Rank Conjecture in Communication Complexity}{Log-Rank Conjecture in Communication Complexity}}\label{144}
\subsection{Definitions and mathematical statement}
For $f:X\times Y\to\{-1,1\}$, Alice knows $x$, Bob knows $y$, and $D(f)$ is the minimum worst-case number of communicated bits in a deterministic protocol computing $f(x,y)$. Its sign matrix is $M_f=(f(x,y))$. The conjecture is
\[
 \exists C\ge1\ \forall f:\qquad
 D(f)\le C\bigl(\log_2\max\{2,\operatorname{rank}_{{\mathbb{R}}}M_f\}\bigr)^C.
\]
The same constant works for all nonempty finite $X,Y$. The rank is over ${\mathbb{R}}$; it is not Boolean rank or nonnegative rank, and the communication is deterministic.
\par\smallskip\textit{Formulation and concept references:} \href{https://arxiv.org/html/1306.1877v2}{[S1]}.

\subsection{Short English statement}
Is deterministic communication complexity always bounded by a fixed power of the logarithm of the sign matrix's real rank?
\subsection{Sources}
\begin{itemize}
\item[S1]S. Lovett, 'Communication is bounded by root of rank', arXiv:1306.1877, 2013; also in Proc. 29th IEEE Conf. on Computational Complexity (CCC), 2014.
\url{https://arxiv.org/html/1306.1877v2}.
\textit{Formulation source.}
\end{itemize}
\end{document}
